% % % % % % % % % % % % % % % % % % % % % % % % % % % % % % % % % % % % % % % % % % % %
%                                                                                     %
% Short Sectioned Assignment LaTeX Template Version 1.0 (5/5/12)                      %
% This template has been downloaded from: http://www.LaTeXTemplates.com               %
%                                                                                     %
% Original author:  Frits Wenneker (http://www.howtotex.com)                          %
%                                                                                     %
% Modified by: Fco Javier Sueza Rodríguez (fcosueza@disroot.org)                      %
%                                                                                     %
% Changes:                                                                            %
%           - Document type scrarticle                                                %
%           - Use babel-lang-spanish package and marvosym                             %
%           - Use hyperref, enumitem and tcolorbox                                    %
%           - Use fancyhdr package an configure it                                    %
%           - Use Time New Roman (mathptmx), Helvetic and Courier fonts               %
%                                                                                     %
% License: CC BY-NC-SA 3.0 (http://creativecommons.org/licenses/by-nc-sa/3.0/)        %
%                                                                                     %
% % % % % % % % % % % % % % % % % % % % % % % % % % % % % % % % % % % % % % % % % % % %

%-----------------------------------------------%
%	              Packages                  %
%-----------------------------------------------%

\documentclass[bibliography=totoc]{scrarticle}

% ---- Text Input/Output ----- %

\usepackage[T1]{fontenc}
\usepackage[utf8]{inputenc}
\usepackage{mathptmx}
\usepackage[scaled=.92]{helvet}
\usepackage{courier}
\usepackage[indent=12pt]{parskip}
\usepackage[titletoc, title]{appendix}

\usepackage{geometry}
\geometry{verbose,tmargin=3cm,bmargin=3cm,lmargin=2.5cm,rmargin=2.5cm}

% ---- Language ----- %

\usepackage[spanish, es-tabla, es-noquoting]{babel}
\usepackage{marvosym}

% ---- Another packages ---- %

\usepackage[toc, sort=use]{glossaries}
\usepackage{amsmath,amsfonts,amsthm}
\usepackage{graphics,graphicx}
\usepackage{tcolorbox}
\usepackage{hyperref}
\usepackage{enumitem}
\usepackage{fancyhdr}
\usepackage{float}
\usepackage{ulem}


%--------------------------------------------------------------------%
%                      Customizing Document                          %
%--------------------------------------------------------------------%

% -------------- Customize headers and footers---------------------- %

\pagestyle{fancy}

\fancyhead[]{}
\fancyfoot[L]{}
\fancyfoot[C]{}
\fancyfoot[R]{\thepage}

\renewcommand{\headrulewidth}{0pt} % Remove header underlines
\renewcommand{\footrulewidth}{0pt} % Remove footer underlines

% Needed to get first page footer correctly
\fancypagestyle{plain}{
    \fancyhf{}
    \fancyfoot[R]{\thepage}
}

% --------------------- Packages Configuration --------------------- %


\hypersetup{
    colorlinks=true,
    linkcolor=magenta,
    urlcolor=magenta
}

\setcounter{secnumdepth}{4}
\setcounter{tocdepth}{4}
\graphicspath{{./images/}}

\setlength{\headheight}{13.6pt} % Customize the height of the header
\numberwithin{figure}{section} % Number figures within sections
\numberwithin{table}{section}

% ------------------------ New Commands ----------------------------- %

\newcommand{\horrule}[1]{\rule{\linewidth}{#1}} % Create horizontal rule command

% ------------------------ Anexx now with name in spanish before title ----------------------- %

\addto\captionsspanish{%
	\renewcommand\appendixname{Anexo}
	\renewcommand\appendixpagename{Anexos}
}




%----------------------------------------------------------------------------------------
%	TÍTULO Y DATOS DEL ALUMNO
%----------------------------------------------------------------------------------------

\title{
\vspace{10ex}
\normalfont \normalsize
\huge \textbf{Amazon Rekognition vs Google Cloud Vision}
}
\author{Francisco Javier Sueza Rodríguez}
\date{\normalsize\today}

\glsenablehyper
\makeglossaries
\loadglsentries{glos}

%----------------------------------------------------------------------------------------
%                                     DOCUMENTO
%----------------------------------------------------------------------------------------
\begin{document}


\maketitle

\thispagestyle{empty}

\vspace{75ex}

\begin{center}
    \begin{tabular}{l l}
        \textbf{Centro}: & IES Zaidin-Vergeles \\
        \textbf{Curso}: & Inteligencia Artificial y Big Data\\
        \textbf{Módulo}: & Programación de Inteligencia Artificial\\
    \end{tabular}
\end{center}

\newpage

\begingroup
\hypersetup{linkcolor=black}
\tableofcontents
\endgroup

\newpage

\begingroup
\hypersetup{linkcolor=black}
\listoffigures
\listoftables
\endgroup

\newpage

\section{Propuesta de Proyecto}
Un \textbf{clan competitivo} del juego Hell Let Loose, un simulador militar de la Segunda Guerra Mundial, esta encontrando \textbf{problemas para desarrollar estrategias efectivas} en sus partidas de liga. Dado que no se encuentra una solución eficiente al problema y resulta bastante complejo el análisis de todo lo que sucede en una partida, se ha decidido crear una aplicación para que \textbf{ayude} en el proceso de \textbf{análisis y toma de decisiones} en la elaboración de la estrategia.

Los \textbf{usuarios} de esta aplicación serían los miembros del clan que han encargado el proyecto, aunque se baraja la posibilidad de que el uso de la aplicación se expanda a otros clanes o miembros de la comunidad competitiva de dicho juego.

En este contexto, la \textbf{inteligencia artificial} nos va a aportar una gran capacidad de análisis, tanto de imágenes como de datos, para ayudar en la elaboración de estrategias efectivas que ayuden al clan a ganar sus partidas, o por lo menos a aumentar el rendimiento en éstas.

\section{Objetivos}
El \textbf{objetivo principal}  de la aplicación es realizar el \textbf{análisis de las partidas}, mediante el \textbf{procesamiento de las capturas de pantalla} proporcionados por lo usuarios, así como de los \textbf{datos generados} durante la partida y que se almacenan en formato CSV. Una vez realizado el análisis, deberá \textbf{ayudar en el desarrollo de estrategias} que se adapten a los diferentes mapas y que ayuden al equipo a ser más eficiente en las partidas. Este servicio se proporcionaría mediante un \textbf{chatbot} al que se puede acceder desde una \textbf{interfaz web}.

Además de este objetivo principal, estaría interesante agregar algunos \textbf{objetivos secundarios}, como pueden ser los siguientes:

\begin{itemize}
	\item \textbf{Integración con Discord}: sería interesante hacer que la aplicación se integrará en el servidor de discord del clan, proporcionando algunas las funcionalidades sin necesidad de acceder a la web.
	\item \textbf{Perfiles de Equipos Rivales}: además de realizar análisis de las estrategias propias, sería bastante útil que pasando capturas de partidas de otros rivales se analizara las estrategias que llevan acabo elaborando un perfil y almacenándolo para su posterior uso o consulta.
	\item \textbf{Perfil Personal}: mediante el análisis de los datos de un usuario concreto, elaborar un perfil individual del jugador,  indicando en que mapas se desempeña mejor, con que armas, en que posiciones, etc, ofreciendo finalmente algunos consejos para mejorar el rendimiento en las partidas y almacenando el resultado, para contrastar con datos futuros y analizar la evolución.
\end{itemize}

\section{Análisis de Riesgos}
En este apartado vamos a realizar un \textbf{análisis de los riesgos} asociados al proyecto. Dada su naturaleza, podemos encontrar menos riesgos que en otro tipos de proyecto, pero sin estar exento de ellos. 

Los \textbf{principales riesgo} asociados a este proyectos son los siguientes: 

\begin{itemize}
	\item \textbf{Riesgos Técnicos}
	\begin{itemize}
		\item \textbf{Precisión en la identificación de elementos}: los modelos podrían no identificar correctamente todos los elementos que aparecen en el mapa, realizando una análisis equivocado lo que afectaría a su función principal. La forma de \textbf{mitigar} este riesgos es proporcionar, durante su entrenamiento, un conjunto de datos e imágenes especificando las características de cada elemento y su uso en el juego.
		
		\item \textbf{Datos de Entrenamiento Insuficientes}: dada las características de la aplicación, recopilar una cantidad de datos suficiente para que los análisis sean precisos, puede ser problemático. Para \textbf{mitigar} este riesgo, se deberá realizar una buena planificación del entrenamiento, proporcionando bastantes capturas de pantalla, limitando inicialmente, si fuera necesario, el alcance a ciertos elementos de los mapas.
		
		\item \textbf{Fiabilidad del análisis}: la distribución de los jugadores y los elementos en el mapa no siempre permiten determinar is una decisión táctica es correcta. Una forma de \textbf{mitigar} este riesgo es mediante la inclusión del contexto de las capturas y la revisión de la conclusión por jugadores experimentados.
		
		\item \textbf{Cantidad de capturas y limitación temporal}: si las capturas de una partida están demasiado espaciadas entre ellas, la evaluación de la situación y de los cambios en la partida pueden no ser correctos. Habrá que \textbf{mitigare} este riesgo proporcionando una cantidad suficiente de capturas, preferiblemente cada pocos minutos, parta que el análisis sea más preciso.
	\end{itemize}
	
	\item \textbf{Riesgos económicos}
	\begin{itemize}
		\item \textbf{Coste de la Infraestructura}: la infraestructura para la ejecución de modelos de IA, especialmente orientados a las visión artificial, puede ser costoso. Una de forma de \textbf{mitigar} este riesgo sería usar modelos ya preentrenados y aprovechar recursos gratuitos, antes de contratar infraestructura en la nube.
		\item \textbf{Coste de Servicios de Inteligencia Artificial}: si se usan APIs comerciales en la aplicación el costo puede aumentar rápidamente con el número de peticiones. Para \textbf{mitigar} es riesgo se pueden establecer límites, explorar las opciones gratuitas o valorar le ejecución de modelos locales o en la nube.
		\item \textbf{Coste de mantenimiento y escalabilidad}: si finalmente se decide ofrecer la aplicación al resto de clanes de la comunidad esto puede tener un incremento sustancial en el coste de mantenimiento del servicios. Se podría \textbf{mitigar} este riesgo planteando la inclusión de otros clanes de la comunidad mediante un servicios de suscripción que cobra los gatos extra generados.
	\end{itemize}
	
	\item \textbf{Riesgos Organizativos}
	\begin{itemize}
		\item \textbf{Planificación Inadecuada}: se puede correr el riesgo de no realizar una planificación adecuada del desarrollo del proyecto. Una forma de \textbf{mitigarlo} es siguiendo un correcto ciclo de desarrollo y planificar correctamente cada fase.
		\item \textbf{Disponibilidad de los Usuarios Colaboradores}: para el correcto entrenamiento y recopilación de datos necesitamos la colaboración de los usuario de la aplicación, para que tomen capturas en situaciones reales de partida. Esto puede resultar en inconvenientes a la hora de organizarse con la gente del clan. Para \textbf{mitigarlo}, deberemos acordar con antelación la colaboración de los miembros del clan, ya sea de todos o de un grupo, para aportar dichas capturas, así como para etiquetar algunos datos y validar los resultados.
		\item \textbf{Alcance Excesivo}: otro riesgo es que queramos que el proyecto haga mas de lo que realmente puede con el presupuesto y la tecnología que se cuenta. Para \textbf{mitigar} es riesgo, podemos establecer límites claros, dejando en primer lugar los objetivos secundarios al margen y centrándonos en la funcionalidad principal.
	\end{itemize}
	
	\item \textbf{Riesgos con los Datos y Calidad de Predicciones}
	\begin{itemize}
		\item \textbf{Recomendaciones tácticas poco efectivas}: puede que el modelo llegue a conclusiones que no son efectivas en las partidas, por diferentes motivos, y que pueden ser contraproducentes para el equipo. Para \textbf{mitigar} estos problemas, se necesitará la validación, por parte de jugadores experimentados, de las estrategias sugeridas, así como proporcionar feedback al sistemas después de aplicar las tácticas sugeridas.
		\item \textbf{Análisis de rivales poco realista}: los análisis de rivales pueden ser poco realistas o equivocados, llevando a una planificación táctica errónea. Para \textbf{mitigar} este riesgo, se deberán proporcionar todas las capturas he información necesaria sobre los rivales para que realice un perfil táctico más ajustado a la realidad.                  
	\end{itemize}
\end{itemize}

Respecto a los \textbf{riesgos éticos}, dada la naturaleza de la aplicación y que esta se va a implementar para un juego online, no se han detectado riesgos de este tipo.

\newpage

\bibliography{citas}
\bibliographystyle{unsrt}


\end{document} 